\section{Моделирование предметной области и разработка требований к программному средству}

Второй раздел посвящен формализации предметной области и определению требований к разрабатываемому программному средству. На основе анализа, выполненного в первом разделе, необходимо описать текущее состояние процесса сквозной аналитики \eng{e-commerce} проекта, выявить проблемы ручной обработки данных и определить целевое состояние, при котором основные операции выполняются в единой информационной системе.

Моделирование предметной области позволяет перейти от общего описания проблемы к конкретным пользовательским сценариям, ролям, функциям и ограничениям системы. Для дипломного проекта это особенно важно, поскольку программное средство должно демонстрировать не отдельную аналитическую панель, а полный цикл сопровождения аналитики: подачу заявки, выбор \eng{SEO}-специалиста, заполнение данных о бизнесе, ввод событий и заказов, проверку данных, расчет показателей, подготовку рекомендаций и публикацию отчета.

При построении модели предметной области необходимо учитывать, что разрабатываемая система не заменяет интернет-магазин и не является рекламным кабинетом. Она располагается на уровне аналитического сопровождения проекта и работает с уже подготовленными сведениями о трафике, событиях, заказах и возвратах. Такой подход позволяет четко определить границы программного средства: оно принимает данные, проверяет их, сохраняет в реляционной базе, рассчитывает показатели и представляет результат пользователям с разными ролями. Внешние источники данных не рассматриваются как отдельные акторы, поскольку в рамках дипломного проекта используется табличный ввод и генерация демонстрационного набора.

\subsection{Анализ и формализация бизнес-процессов предметной области}

В текущем состоянии процесс анализа эффективности \eng{e-commerce} проекта часто выполняется вручную. Клиент передает данные о сайте, заказах, возвратах и источниках трафика в виде отдельных таблиц или файлов. \eng{SEO}-специалист проверяет их структуру, сопоставляет переходы с заказами, рассчитывает показатели и формирует отчет. Такой подход не обеспечивает единого хранилища данных, затрудняет повторное использование результатов и увеличивает вероятность ошибок.

Еще одной особенностью ручного процесса является отсутствие единого жизненного цикла заявки. До внедрения программного средства обращение клиента может существовать в переписке, таблице, личных заметках специалиста или отдельном документе. При этом трудно определить, когда заявка была принята, какие данные уже получены, требуется ли уточнение и опубликован ли отчет. Система должна устранить этот разрыв за счет статусов, связанных сущностей и централизованного хранения истории действий.

Модель процесса \eng{AS-IS} отражает ручной порядок выполнения аналитических операций до внедрения программного средства. В этом варианте большая часть работы выполняется специалистом вручную, а клиент ожидает итоговый отчет без прозрачного доступа к промежуточным результатам. Модель \eng{AS-IS} процесса анализа эффективности \eng{e-commerce} представлена на рисунке~\ref{fig:as-is-ecommerce}.

\begin{figure}[H]
    \centering
    \diplomaimage{images/fig-2-1-as-is-ecommerce-analytics.png}{80mm}
    \caption{Модель \eng{AS-IS} процесса анализа эффективности \eng{e-commerce}}
    \label{fig:as-is-ecommerce}
\end{figure}

В модели \eng{AS-IS} можно выделить несколько ключевых этапов: обращение клиента, описание проблемы, передача таблиц, ручная проверка данных, расчет показателей, подготовка отчета и отправка результата клиенту. Наиболее трудоемкими являются операции проверки данных, сопоставления источников трафика с заказами и подготовки итогового отчета. Длительность выполнения функций процесса приведена в таблице~\ref{tab:analytics-duration}.

\begin{table}[H]
    \caption{Длительность выполнения функций процесса}
    \label{tab:analytics-duration}
    {\small
    \setlength{\tabcolsep}{4pt}
    \begin{tabularx}{\textwidth}{|P{0.32\textwidth}|Y|Y|Y|}
        \hline
        Функция процесса & Ручное выполнение, мин & Выполнение в системе, мин & Сокращение времени \\
        \hline
        Сбор данных о переходах, событиях, заказах и возвратах & 120 & 35 & 85 \\
        \hline
        Проверка структуры и полноты табличных данных & 90 & 20 & 70 \\
        \hline
        Сопоставление источников трафика с заказами & 110 & 25 & 85 \\
        \hline
        Расчет показателей эффективности проекта & 75 & 10 & 65 \\
        \hline
        Подготовка аналитического отчета и рекомендаций & 100 & 30 & 70 \\
        \hline
        Итого & 495 & 120 & 375 \\
        \hline
    \end{tabularx}
    }
\end{table}

Сравнение времени показывает, что наибольший эффект автоматизации достигается за счет сокращения ручных операций, связанных с проверкой и сопоставлением данных. При использовании программного средства специалист продолжает выполнять аналитическую интерпретацию результата, но рутинные операции расчета и подготовки данных выполняются системой. Гистограмма анализа времени выполнения аналитических операций представлена на рисунке~\ref{fig:analytics-time-histogram}.

При этом автоматизация не должна подменять аналитическое мышление специалиста. Система может показать, что определенный источник имеет низкую конверсию или высокую долю возвратов, но вывод о причинах проблемы формулирует \eng{SEO}-специалист. Именно поэтому в модели \eng{TO-BE} сохраняется этап анализа результатов и подготовки рекомендаций. Программное средство выступает инструментом поддержки решения, а не полностью автономным экспертным модулем.

\begin{figure}[H]
    \centering
    \diplomaimage{images/fig-2-2-analytics-time-histogram.png}{80mm}
    \caption{Гистограмма анализа времени выполнения аналитических операций}
    \label{fig:analytics-time-histogram}
\end{figure}

Ручная организация процесса приводит не только к увеличению времени работы, но и к возникновению ошибок. Проблемы исполнения процесса связаны с отсутствием единой базы данных, разрозненностью файлов, отсутствием автоматической проверки формата и невозможностью быстро обновить показатели при изменении исходных данных. Описание проблем исполнения процесса приведено в таблице~\ref{tab:analytics-problems}.

\begin{table}[H]
    \caption{Описание проблем исполнения процесса}
    \label{tab:analytics-problems}
    {\small
    \setlength{\tabcolsep}{4pt}
    \begin{tabularx}{\textwidth}{|P{0.28\textwidth}|L|L|}
        \hline
        Проблема & Проявление в процессе & Причина возникновения \\
        \hline
        Разрозненность данных & Переходы, события, заказы и возвраты хранятся в разных таблицах & Отсутствие единого программного средства для сопровождения аналитики \\
        \hline
        Высокая доля ручных операций & Специалист вручную проверяет строки, периоды, источники и суммы заказов & Нет автоматизированной проверки структуры и полноты данных \\
        \hline
        Сложность повторного расчета & При изменении исходных данных отчет приходится пересобирать заново & Показатели не связаны с централизованной моделью данных \\
        \hline
        Низкая прозрачность для клиента & Клиент видит только готовый отчет и не отслеживает состояние заявки & Нет единого интерфейса для заявок, проектов и отчетов \\
        \hline
        Риск методических ошибок & Источники трафика и заказы могут сопоставляться по разным правилам & Правила аналитики фиксируются не в системе, а в рабочих файлах специалиста \\
        \hline
    \end{tabularx}
    }
\end{table}

Каждая из выявленных проблем приводит к потерям, рискам и последствиям для участников процесса. Для клиента это выражается в задержке получения отчета и неполной прозрачности результатов. Для \eng{SEO}-специалиста -- в росте ручной нагрузки и необходимости повторять однотипные операции. Для администратора -- в сложности контроля пользователей, проектов и демонстрационных данных. Потери, риски и последствия при выполнении процесса представлены в таблице~\ref{tab:analytics-risks}.

\begin{table}[H]
    \caption{Потери, риски и последствия при выполнении процесса}
    \label{tab:analytics-risks}
    {\small
    \setlength{\tabcolsep}{4pt}
    \begin{tabularx}{\textwidth}{|P{0.24\textwidth}|L|L|L|}
        \hline
        Фактор & Потери & Риск & Последствие \\
        \hline
        Ручная проверка данных & Увеличение времени подготовки аналитики & Пропуск некорректной строки или дублирование записи & Искажение показателей конверсии, выручки и возвратов \\
        \hline
        Отдельные файлы по проектам & Сложность поиска актуальной версии данных & Использование устаревшей таблицы & Подготовка отчета на основании неполных сведений \\
        \hline
        Ручной расчет показателей & Дополнительная нагрузка на специалиста & Ошибка в формуле или периоде анализа & Неверные выводы об эффективности источников трафика \\
        \hline
        Отсутствие единого статуса заявки & Необходимость уточнять состояние работы вручную & Потеря контроля над этапом обработки & Увеличение срока подготовки отчета \\
        \hline
    \end{tabularx}
    }
\end{table}

Для устранения выявленных проблем необходимо автоматизировать основные этапы процесса. Автоматизация не должна исключать аналитическую роль специалиста: система рассчитывает показатели и формирует основу для отчета, а \eng{SEO}-специалист интерпретирует результаты и готовит рекомендации. Мероприятия, направленные на автоматизацию процессов, приведены в таблице~\ref{tab:analytics-automation-measures}.

\begin{table}[H]
    \caption{Мероприятия, направленные на автоматизацию процессов}
    \label{tab:analytics-automation-measures}
    {\small
    \setlength{\tabcolsep}{4pt}
    \begin{tabularx}{\textwidth}{|P{0.27\textwidth}|L|L|}
        \hline
        Мероприятие & Содержание мероприятия & Ожидаемый результат \\
        \hline
        Регистрация заявок в системе & Клиент создает заявку, выбирает специалиста и описывает цель аналитики & Централизованный учет обращений и статусов \\
        \hline
        Ввод данных в табличном интерфейсе & Клиент передает переходы, события, заказы, возвраты и расходы через формы системы & Снижение зависимости от внешних файлов \\
        \hline
        Автоматическая проверка данных & Система проверяет обязательные поля, форматы дат, суммы и связи записей & Уменьшение количества ошибок в расчетах \\
        \hline
        Расчет аналитических показателей & Система рассчитывает переходы, заказы, выручку, конверсию, средний чек и возвраты & Быстрое получение показателей по проекту \\
        \hline
        Формирование отчетов и рекомендаций & Специалист подготавливает выводы на основе рассчитанных показателей & Повышение прозрачности результатов для клиента \\
        \hline
    \end{tabularx}
    }
\end{table}

Целевое состояние процесса описывается моделью \eng{TO-BE}. В этой модели клиент подает заявку в системе, выбирает \eng{SEO}-специалиста, заполняет данные о бизнесе и передает табличные сведения. Специалист принимает заявку, проводит диагностику, настраивает модель аналитики, после чего система проверяет данные, рассчитывает показатели и формирует основу для \eng{dashboard} и отчета. Модель \eng{TO-BE} процесса сквозной аналитики \eng{e-commerce} представлена на рисунке~\ref{fig:to-be-ecommerce}.

В модели \eng{TO-BE} данные становятся центральным объектом процесса. После создания проекта все сведения связываются с ним через единую модель: источники трафика относятся к проекту, события и заказы фиксируются по проекту и источнику, возвраты связываются с заказами, отчеты и рекомендации создаются по результатам анализа. Благодаря этому один и тот же набор данных может использоваться для расчета \eng{KPI}, построения воронки, анализа источников и подготовки отчета. Такая связность является основным преимуществом целевой модели по сравнению с ручным процессом.

Целевая модель также учитывает необходимость демонстрации системы в учебных условиях. Если клиент не вводит данные вручную, администратор может сгенерировать демонстрационный набор для выбранного проекта. Это позволяет показать расчет показателей, работу \eng{dashboard}, формирование рекомендаций и публикацию отчета даже при отсутствии внешнего интернет-магазина. Такой механизм не заменяет реальные интеграции, но соответствует границам дипломного проекта и делает систему воспроизводимой.

\begin{figure}[H]
    \centering
    \diplomaimage{images/fig-2-3-to-be-ecommerce-analytics.png}{80mm}
    \caption{Модель \eng{TO-BE} процесса сквозной аналитики \eng{e-commerce}}
    \label{fig:to-be-ecommerce}
\end{figure}

Автоматизация бизнес-процессов позволяет связать организационные действия пользователей с функциями программного средства. Процесс подачи заявки связывается с модулем заявок, ввод данных -- с табличным интерфейсом, расчет показателей -- с аналитическим модулем, а публикация результата -- с отчетами и рекомендациями. Автоматизация бизнес-процессов приведена в таблице~\ref{tab:analytics-process-automation}.

\begin{table}[H]
    \caption{Автоматизация бизнес-процессов}
    \label{tab:analytics-process-automation}
    {\small
    \setlength{\tabcolsep}{4pt}
    \begin{tabularx}{\textwidth}{|P{0.26\textwidth}|L|L|}
        \hline
        Бизнес-процесс & Реализация в программном средстве & Результат автоматизации \\
        \hline
        Подключение аналитики проекта & Создание заявки, выбор \eng{SEO}-специалиста, заполнение брифа & Заявка получает статус и ответственного исполнителя \\
        \hline
        Сбор исходных данных & Ввод переходов, событий, заказов, возвратов и расходов через табличные формы & Данные сохраняются в единой базе проекта \\
        \hline
        Проверка качества данных & Валидация обязательных полей, типов, связей и периодов & Пользователь получает список ошибок или подтверждение корректности \\
        \hline
        Расчет аналитики & Автоматический расчет \eng{KPI}, воронки продаж и эффективности источников & Показатели становятся доступными на \eng{dashboard} \\
        \hline
        Подготовка результата & Создание рекомендаций, публикация отчета, просмотр клиентом & Итог аналитической работы фиксируется в системе \\
        \hline
    \end{tabularx}
    }
\end{table}

Таким образом, переход от модели \eng{AS-IS} к модели \eng{TO-BE} направлен на уменьшение ручного труда, повышение прозрачности статусов и централизованное хранение данных. Предложенная автоматизация сохраняет участие специалиста в аналитической интерпретации, но переносит повторяемые операции проверки, расчета и подготовки данных в программное средство.

Результатом формализации бизнес-процессов является набор функций, которые должны быть отражены в требованиях. Система должна поддерживать учет пользователей, заявок и проектов, прием табличных данных, проверку качества сведений, расчет аналитических показателей, отображение результатов и подготовку отчетов. Кроме того, необходимо предусмотреть административные функции, поскольку без управления пользователями, ролями и демонстрационными данными система не будет завершенной с точки зрения учебного программного средства.

\subsection{Разработка спецификации требований к программному средству}

На основании формализованного процесса необходимо определить действующих лиц, варианты использования и требования к программному средству. В системе выделяются три роли: клиент, \eng{SEO}-специалист и администратор. Внешнее \eng{e-commerce} приложение не включается в список акторов, поскольку в дипломном проекте данные поступают через табличный ввод пользователя или через генерацию демонстрационного массива администратором.

Выделение именно трех ролей позволяет сохранить систему понятной и одновременно отразить основные границы ответственности. Клиент отвечает за постановку задачи и передачу исходных данных. \eng{SEO}-специалист отвечает за методическую часть аналитики, проверку данных, настройку модели и интерпретацию результатов. Администратор отвечает за системное сопровождение, учетные записи, демонстрационные данные и служебные операции. Если добавить внешнее приложение как отдельного участника, границы проекта расширятся до разработки интеграций, что не соответствует учебной постановке задачи.

При описании требований необходимо учитывать, что одна и та же сущность может быть доступна разным ролям, но с разными правами. Например, клиент может просматривать только свои проекты и отчеты, \eng{SEO}-специалист -- назначенные ему заявки и проекты, а администратор -- системные справочники и учетные записи. Поэтому требования к функциям системы должны рассматриваться вместе с требованиями к разграничению доступа. В дальнейшем это будет отражено в архитектуре серверной части и механизмах авторизации.

Диаграмма вариантов использования фиксирует границы системы и основные функции пользователей. Клиент работает с заявками, брифом, исходными данными, \eng{dashboard} и отчетами. \eng{SEO}-специалист принимает заявки, проводит диагностику, настраивает модель аналитики, анализирует показатели и формирует отчет. Администратор управляет пользователями, проектами, справочниками, резервным копированием и демонстрационными данными. Диаграмма вариантов использования системы сквозной аналитики представлена на рисунке~\ref{fig:use-case-analytics}.

Варианты использования сгруппированы вокруг жизненного цикла проекта. Сначала клиент инициирует работу через заявку и бриф. Затем специалист принимает заявку, выполняет диагностику и настраивает правила анализа. После передачи данных система выполняет проверку и расчет показателей. На заключительном этапе специалист формирует рекомендации и отчет, а клиент просматривает результат. Администратор обеспечивает системные условия для выполнения этого цикла: управляет пользователями, ролями, проектами и демонстрационными данными.

\begin{figure}[H]
    \centering
    \diplomaimage{images/fig-2-4-use-case-analytics-system.png}{82mm}
    \caption{Диаграмма вариантов использования системы сквозной аналитики}
    \label{fig:use-case-analytics}
\end{figure}

Действующие лица системы описаны в таблице~\ref{tab:analytics-actors}.

\begin{table}[H]
    \caption{Описание действующих лиц}
    \label{tab:analytics-actors}
    {\small
    \setlength{\tabcolsep}{4pt}
    \begin{tabularx}{\textwidth}{|P{0.25\textwidth}|L|L|}
        \hline
        Действующее лицо & Назначение & Основные действия \\
        \hline
        Клиент & Представитель \eng{e-commerce} проекта, передающий исходные данные и получающий аналитический результат & Создает заявку, выбирает специалиста, заполняет бриф, вводит данные, просматривает \eng{dashboard} и отчет \\
        \hline
        \eng{SEO}-специалист & Основной рабочий пользователь, отвечающий за диагностику проекта и подготовку аналитических выводов & Принимает заявку, проверяет данные, настраивает модель аналитики, анализирует показатели, формирует рекомендации и отчет \\
        \hline
        Администратор & Системный пользователь, обеспечивающий управление учетными записями и служебными данными & Управляет пользователями, ролями, проектами, справочниками, резервными копиями и демонстрационными данными \\
        \hline
    \end{tabularx}
    }
\end{table}

Описание вариантов использования приведено в таблице~\ref{tab:analytics-use-cases}. Варианты использования охватывают полный цикл работы системы: от подачи заявки до формирования отчета и управления системными данными.

{\small
\setlength{\tabcolsep}{4pt}
\begin{longtable}{|P{0.11\textwidth}|P{0.27\textwidth}|P{0.18\textwidth}|P{0.34\textwidth}|}
    \caption{Описание вариантов использования}
    \label{tab:analytics-use-cases}\\
    \hline
    Код & Название & Роль & Краткое описание \\
    \hline
    \endfirsthead
    \caption*{Продолжение таблицы \thetable}\\
    \hline
    Код & Название & Роль & Краткое описание \\
    \hline
    \endhead
    \eng{UC-1} & Подать заявку на подключение сквозной аналитики & Клиент & Клиент указывает сайт, нишу, проблему, цель аналитики и отправляет заявку в систему \\
    \hline
    \eng{UC-2} & Выбрать \eng{SEO}-специалиста & Клиент & Клиент выбирает специалиста из списка активных пользователей системы \\
    \hline
    \eng{UC-3} & Заполнить данные о бизнесе & Клиент & Клиент заполняет бриф: регион, тип товаров, каналы продаж, средний чек и период анализа \\
    \hline
    \eng{UC-4} & Передать данные о продажах и событиях & Клиент & Клиент вводит переходы, события, покупки, возвраты и источники трафика в табличном виде \\
    \hline
    \eng{UC-5} & Просмотреть \eng{dashboard} и отчет & Клиент & Клиент просматривает рассчитанные показатели, воронку продаж, рекомендации и опубликованный отчет \\
    \hline
    \eng{UC-6} & Принять заявку в работу & \eng{SEO}-специалист & Специалист открывает назначенную заявку, проверяет первичные данные и меняет статус на обработку \\
    \hline
    \eng{UC-7} & Провести диагностику \eng{e-commerce} проекта & \eng{SEO}-специалист & Специалист анализирует сайт, цели, источники трафика и корректность исходных данных \\
    \hline
    \eng{UC-8} & Настроить модель аналитики проекта & \eng{SEO}-специалист & Специалист задает источники трафика, целевые события, этапы воронки и правила сопоставления данных \\
    \hline
    \eng{UC-9} & Проанализировать эффективность источников и воронку продаж & \eng{SEO}-специалист & Система рассчитывает показатели, а специалист оценивает конверсию, выручку, возвраты и слабые этапы воронки \\
    \hline
    \eng{UC-10} & Сформировать рекомендации и отчет & \eng{SEO}-специалист & Специалист подготавливает выводы, рекомендации и публикует отчет для клиента \\
    \hline
    \eng{UC-11} & Управлять пользователями и ролями & Администратор & Администратор создает учетные записи, назначает роли, блокирует или активирует пользователей \\
    \hline
    \eng{UC-12} & Управлять системными данными проекта & Администратор & Администратор создает проекты вручную, генерирует демонстрационные данные, управляет справочниками, резервным копированием и журналом действий \\
    \hline
\end{longtable}
}

Между вариантами использования существуют связи включения. Подача заявки включает выбор \eng{SEO}-специалиста и заполнение данных о бизнесе. Анализ эффективности источников использует предварительно настроенную модель аналитики. Формирование отчета включает анализ показателей и подготовку рекомендаций. Управление системными данными может включать генерацию демонстрационного массива для проекта.

Функциональные требования определяют поведение системы и перечень функций, которые должны быть реализованы. Они сформированы на основе ролей пользователей, вариантов использования и целевого процесса \eng{TO-BE}. Функциональные требования к системе сквозной аналитики приведены в таблице~\ref{tab:functional-requirements}.

Функциональные требования связаны с предметной областью и должны быть проверяемыми. Если требование формулируется как возможность создания заявки, то в системе должен существовать интерфейс для ввода данных заявки, серверный обработчик для сохранения данных и запись в базе. Если требование описывает расчет показателей, то система должна иметь набор исходных данных, правила расчета и механизм отображения результата. Такой подход позволяет в дальнейшем связать требования с архитектурой, моделью данных и тестовыми сценариями.

Отдельное внимание уделяется требованиям к данным. Для расчета показателей недостаточно хранить только итоговую сумму выручки. Система должна иметь сведения о переходах, событиях, заказах, возвратах и источниках трафика. Это позволяет рассчитывать показатели в разных разрезах: по проекту, периоду, источнику и этапу воронки. Поэтому требования к вводу и проверке данных являются базовыми для всей системы.

{\small
\setlength{\tabcolsep}{4pt}
\begin{longtable}{|P{0.12\textwidth}|P{0.40\textwidth}|P{0.38\textwidth}|}
    \caption{Функциональные требования к системе сквозной аналитики}
    \label{tab:functional-requirements}\\
    \hline
    Код & Требование & Обоснование \\
    \hline
    \endfirsthead
    \caption*{Продолжение таблицы \thetable}\\
    \hline
    Код & Требование & Обоснование \\
    \hline
    \endhead
    \eng{FR-1} & Система должна обеспечивать авторизацию пользователей & Доступ к заявкам, проектам и отчетам должен предоставляться только зарегистрированным пользователям \\
    \hline
    \eng{FR-2} & Система должна поддерживать роли клиента, \eng{SEO}-специалиста и администратора & Разные пользователи выполняют разные действия и имеют разные права доступа \\
    \hline
    \eng{FR-3} & Система должна позволять клиенту создавать заявку на подключение аналитики & Заявка является начальной точкой процесса сопровождения проекта \\
    \hline
    \eng{FR-4} & Система должна позволять клиенту выбирать \eng{SEO}-специалиста & Заявка должна быть связана с ответственным исполнителем \\
    \hline
    \eng{FR-5} & Система должна позволять заполнять данные о бизнесе & Для диагностики проекта необходимы сведения о сайте, нише, регионе и целях анализа \\
    \hline
    \eng{FR-6} & Система должна поддерживать ввод переходов, событий, покупок и возвратов & Эти данные являются основой расчета показателей сквозной аналитики \\
    \hline
    \eng{FR-7} & Система должна проверять корректность введенных данных & Проверка снижает риск расчета показателей на ошибочных или неполных данных \\
    \hline
    \eng{FR-8} & Система должна рассчитывать аналитические показатели проекта & Пользователю необходимы переходы, заказы, выручка, конверсия, средний чек и возвраты \\
    \hline
    \eng{FR-9} & Система должна отображать \eng{dashboard} проекта & Аналитический результат должен быть представлен в наглядной форме \\
    \hline
    \eng{FR-10} & Система должна поддерживать формирование отчетов и рекомендаций & Итог работы специалиста должен быть зафиксирован и доступен клиенту \\
    \hline
    \eng{FR-11} & Система должна позволять администратору управлять пользователями и ролями & Администратор отвечает за учетные записи и разграничение доступа \\
    \hline
    \eng{FR-12} & Система должна поддерживать генерацию демонстрационных данных & Учебный проект должен демонстрироваться без внешнего интернет-магазина и стороннего \eng{API} \\
    \hline
    \eng{FR-13} & Система должна поддерживать резервное копирование и журнал действий & Необходимо обеспечить восстановление данных и контроль действий пользователей \\
    \hline
\end{longtable}
}

Нефункциональные требования задают ограничения по качеству, производительности, безопасности и условиям эксплуатации. Они не описывают отдельные пользовательские функции, но определяют свойства, которыми должна обладать система. Нефункциональные требования приведены в таблице~\ref{tab:nonfunctional-requirements}.

Выбор технологического стека также отражается в нефункциональных требованиях. Серверная часть реализуется на \eng{Go 1.23}, маршрутизация запросов выполняется через \eng{chi}, работа с базой данных -- через \eng{pgx}, а слой типобезопасного доступа к данным формируется с использованием \eng{sqlc}. Клиентская часть реализуется на \eng{Next.js 15}, \eng{React 19} и \eng{TypeScript}. Для хранения данных применяется \eng{PostgreSQL 16}, развертывание выполняется через \eng{Docker Compose}, а \eng{Caddy} используется для \eng{HTTPS} и обратного прокси.

Нефункциональные требования необходимы для того, чтобы система была не только функционально полной, но и пригодной для демонстрации. Например, требование запуска через \eng{Docker Compose} позволяет воспроизвести систему без ручной настройки окружения. Требование к \eng{HTTPS} показывает учет защиты передаваемых данных. Требование к ролевому доступу подтверждает, что аналитические сведения разных клиентов не должны смешиваться и отображаться посторонним пользователям.

{\small
\setlength{\tabcolsep}{4pt}
\begin{longtable}{|P{0.12\textwidth}|P{0.40\textwidth}|P{0.38\textwidth}|}
    \caption{Нефункциональные требования}
    \label{tab:nonfunctional-requirements}\\
    \hline
    Код & Требование & Обоснование \\
    \hline
    \endfirsthead
    \caption*{Продолжение таблицы \thetable}\\
    \hline
    Код & Требование & Обоснование \\
    \hline
    \endhead
    \eng{NFR-1} & Время отклика основных запросов \eng{API} должно быть не более 500 мс на тестовом наборе данных & Пользовательский интерфейс должен оставаться удобным при работе с заявками, проектами и отчетами \\
    \hline
    \eng{NFR-2} & Страница \eng{dashboard} должна открываться не более чем за 2 секунды на тестовом наборе данных & Основной аналитический экран должен быстро отображать показатели проекта \\
    \hline
    \eng{NFR-3} & Пароли пользователей должны храниться в хешированном виде & Требуется защита учетных данных пользователей \\
    \hline
    \eng{NFR-4} & Передача данных должна выполняться с использованием \eng{HTTPS} при развертывании через \eng{Caddy} & Необходимо защищать данные авторизации и аналитические сведения проекта \\
    \hline
    \eng{NFR-5} & Интерфейс должен работать в актуальных версиях современных браузеров & Пользователи должны иметь доступ к системе без установки специализированного программного обеспечения \\
    \hline
    \eng{NFR-6} & Система должна запускаться через \eng{Docker Compose} без использования среды разработки & Демонстрация дипломного проекта должна быть воспроизводимой \\
    \hline
    \eng{NFR-7} & Данные должны храниться в \eng{PostgreSQL 16} & Реляционная база данных обеспечивает надежное хранение связанных сущностей проекта \\
    \hline
    \eng{NFR-8} & Серверная часть должна быть реализована на \eng{Go 1.23} с использованием \eng{chi}, \eng{pgx} и \eng{sqlc} & Выбранный стек обеспечивает явную структуру \eng{REST API} и типобезопасную работу с базой данных \\
    \hline
    \eng{NFR-9} & Клиентская часть должна быть реализована на \eng{Next.js 15}, \eng{React 19} и \eng{TypeScript} & Интерфейс должен поддерживать современную компонентную разработку и статическую типизацию \\
    \hline
    \eng{NFR-10} & Доступ к данным должен ограничиваться ролью пользователя & Клиент видит только свои проекты и отчеты, специалист -- назначенные проекты, администратор -- системные данные \\
    \hline
\end{longtable}
}

Сформулированные требования задают основу для последующего проектирования архитектуры, модели данных и пользовательского интерфейса. В дальнейшем они будут использованы при описании компонентов системы, структуры базы данных, прикладной логики и механизмов информационной безопасности.

На следующем этапе требования должны быть преобразованы в проектные решения. Функции заявок и проектов будут отражены в соответствующих сущностях базы данных и серверных обработчиках. Требования к расчету показателей определят состав аналитического сервиса. Требования к пользовательским ролям будут использованы при описании механизмов авторизации и проверки прав. Требования к интерфейсу определят набор экранов клиентской части: заявки, проекты, \eng{dashboard}, источники трафика, отчеты, рекомендации и административные страницы.

\subsection{Образ предлагаемого решения}

Предлагаемое решение представляет собой учебное программное средство для автоматизации сквозной аналитики \eng{e-commerce} проектов. Система должна объединять функции регистрации заявок, ведения проектов, ввода аналитических данных, проверки корректности сведений, расчета показателей, формирования рекомендаций и подготовки отчетов. Основная цель решения заключается в том, чтобы показать полный цикл работы с аналитическими данными без разработки внешнего интернет-магазина или подключения стороннего сервиса.

С точки зрения пользователя система должна восприниматься как единая рабочая среда. Клиент начинает работу с заявки, затем передает сведения о бизнесе и данные о событиях, а после обработки получает доступ к аналитическому результату. \eng{SEO}-специалист работает с назначенными заявками и проектами, видит качество данных, выполняет диагностику и формирует выводы. Администратор поддерживает работоспособность системы, управляет ролями и обеспечивает наличие тестовых данных для демонстрации. Такой образ решения соответствует практической логике работы аналитического отдела, но остается выполнимым в рамках дипломного проекта.

Важной частью образа решения является разделение ответственности между клиентской и серверной частями. Клиентская часть отвечает за удобное представление данных, формы ввода, таблицы, фильтры, графики и страницы отчетов. Серверная часть принимает запросы, проверяет права пользователя, обращается к базе данных, рассчитывает показатели и возвращает результат в формате \eng{JSON}. Такое разделение упрощает дальнейшее проектирование и соответствует выбранному стеку: \eng{Next.js 15}, \eng{React 19} и \eng{TypeScript} на клиенте, \eng{Go 1.23}, \eng{chi}, \eng{pgx} и \eng{sqlc} на сервере.

В системе предполагается три основных пользовательских контура. Клиентский контур обеспечивает создание заявки, выбор \eng{SEO}-специалиста, заполнение данных о бизнесе, ввод событий и продаж, просмотр \eng{dashboard} и отчета. Контур \eng{SEO}-специалиста включает прием заявки, диагностику проекта, настройку модели аналитики, проверку данных, анализ источников трафика, формирование рекомендаций и публикацию отчета. Административный контур отвечает за пользователей, роли, проекты, справочники, резервное копирование, журнал действий и генерацию демонстрационных данных.

Ключевой особенностью предлагаемого решения является автономность демонстрации. Система не использует внешнее \eng{e-commerce} приложение как отдельного актора. Данные поступают двумя способами: клиент вводит сведения вручную через табличные формы, а администратор может сформировать демонстрационный массив для выбранного проекта. Такой подход позволяет проверить все основные функции программного средства в рамках дипломного проекта и не зависеть от сторонних учетных записей или внешних интеграций.

Демонстрационный массив данных должен быть достаточно реалистичным, чтобы на его основе можно было показать работу аналитики. В него могут входить источники трафика, переходы за выбранный период, события просмотра товара, добавления в корзину и оформления заказа, оплаченные заказы, возвраты и расходы по каналам. После генерации таких данных система может рассчитать переходы, заказы, выручку, конверсию, средний чек, возвраты, эффективность источников и показатели воронки. Это позволяет продемонстрировать полный сценарий даже без подключения внешних сервисов.

Для обеспечения воспроизводимости решения развертывание должно выполняться через \eng{Docker Compose}. В составе окружения предусматриваются контейнер серверной части, контейнер клиентской части, база данных \eng{PostgreSQL 16} и \eng{Caddy} как обратный прокси. Такое решение упрощает запуск системы и делает его независимым от локальной среды разработки. Для дипломного проекта это важно, поскольку проверяющий должен видеть работоспособность программного средства, а не набор фрагментов кода.

Целевой результат работы системы -- аналитический \eng{dashboard} и отчет, позволяющие оценить эффективность источников трафика, конверсию, выручку, возвраты и слабые этапы воронки продаж. Для клиента результат выражается в прозрачном представлении состояния проекта, для \eng{SEO}-специалиста -- в сокращении ручного труда и систематизации аналитических операций, для администратора -- в контроле пользователей и демонстрационных данных.

\eng{Dashboard} должен содержать показатели, которые непосредственно связаны с задачами сквозной аналитики: количество переходов, количество заказов, выручку, конверсию, средний чек, сумму возвратов и эффективность источников трафика. Отчет должен дополнять визуальные показатели текстовыми выводами и рекомендациями. Важно, чтобы рекомендации не выглядели как рекламные обещания, а были связаны с данными проекта: например, с низкой конверсией определенного источника, высокой долей возвратов или потерями на этапе оформления заказа.

Предлагаемое решение должно учитывать образовательный характер проекта. Оно не описывается как действующий коммерческий сервис и не предполагает фактических продаж. Его назначение состоит в демонстрации процесса разработки и проектирования программного средства: от анализа предметной области и требований до архитектуры, интерфейса, базы данных и механизмов безопасности. Поэтому в дальнейшем экономический раздел должен рассматривать условную эффективность автоматизации, а не реальные коммерческие показатели продукта.

Таким образом, образ предлагаемого решения соответствует выявленным проблемам предметной области и сформулированным требованиям. Разрабатываемое программное средство должно быть не коммерческим продуктом, а учебной \eng{web}-системой, демонстрирующей принципы сквозной аналитики, работу с ролями пользователей, расчет показателей и подготовку отчетности на основе табличных данных.
